% ---------------------------------------------------------------------------
% Corpus construction + sink-family coverage.
% Every number here is traceable:
%   funnel         -> ToFu_src/nazmul.todo, "Funnel on the real data"
%   per-language   -> Benchmarks/extraction/targets.jsonl (1,847 rows)
%   tools + sinks  -> Benchmarks/extraction/results.jsonl (1,778 rows)
%   failures       -> Benchmarks/extraction/failures.jsonl (69 rows)
% Recomputed 2026-09-11. Do not edit a number here without rerunning the source.
% ---------------------------------------------------------------------------

\subsection{Corpus construction}
\label{sec:corpus}

We draw our population from \textsc{MCPCorpus}~\cite{lin2025mcpcorpus}, a
crawl of the \texttt{mcp.so} directory. We use the published
\texttt{mcpso\_servers\_cleaned.json} snapshot at a pinned commit: the
upstream repository has not been updated since 2025-09-01, so the population
is effectively frozen, which we treat as a benefit for reproducibility rather
than a limitation.

Construction proceeds in four stages (Table~\ref{tab:funnel}). From 13,875
records we retain those that name a GitHub repository, have at least ten
stars, and are neither archived nor an MCP \emph{client}, leaving 2,501
records. These records are then deduplicated \emph{by repository}: 279
records collapse, because large monorepos are listed once per server they
contain (\texttt{modelcontextprotocol/servers} appears 35 times,
\texttt{baidu-maps/mcp} 8 times). Skipping this step would let a single
repository dominate any sample. That leaves 2,222 unique repositories.

ToFu ships two runners, for Python and for Node/TypeScript, so we restrict to
those two languages: 808 Python and 1,039 Node, or 1,847 repositories. The
remaining 375 (Go 154, Rust 41, Java 34, C\# 19, PHP 17, and a tail) are a
stated exclusion.

Finally we clone each of the 1,847 repositories at depth one and extract its
tool registrations and sink usage statically. Extraction succeeded for 1,778;
the 69 failures are repositories that have since been deleted or made private
(52), other clone errors (13), and clone timeouts (4). We extract tool
metadata ourselves rather than reading the corpus \texttt{tools} field,
because only 55 of the 13,875 records carry a parsable tool schema.

Of the 1,778 repositories scanned, 1,404 expose at least one tool---597
Python and 807 Node---for a total of 24,194 tools. The remaining 374 register
no tool our extractor recognises, and we exclude them. The two language
populations are strikingly similar in shape: both have a median of 6 tools per
repository (mean 17.0 Python, 17.4 Node) and a long tail reaching 881 and 918
tools respectively, so language is not a confound for tool count.

\begin{table}[t]
\centering
\caption{Corpus construction. Counts are repositories except where noted.}
\label{tab:funnel}
\small
\begin{tabular}{@{}lr@{}}
\toprule
\textbf{Stage} & \textbf{Count} \\
\midrule
\textsc{MCPCorpus} records                            & 13,875 \\
\quad has repo, $\geq$10 stars, not archived, not client & 2,501 \\
\quad deduplicated by repository                      & 2,222 \\
\quad Python or Node/TS (ToFu has two runners)        & 1,847 \\
\quad static extraction succeeded                     & 1,778 \\
\quad exposes $\geq$1 tool                            & \textbf{1,404} \\
\midrule
\quad\quad Python & 597 \\
\quad\quad Node/TS & 807 \\
\midrule
Tools extracted & \textbf{24,194} \\
\quad\quad Python & 10,130 \\
\quad\quad Node/TS & 14,064 \\
\bottomrule
\end{tabular}
\end{table}


\begin{table}[t]
\centering
\caption{Corpus coverage of ToFu's 21 introspection tracers. Each row is the
  fraction of repositories containing at least one static reference to an API
  that the corresponding tracer instruments; a repository may match several
  families, so columns do not sum to 100\%. This measures how much of the
  corpus each tracer can \emph{reach}, not how many defects exist.
  $n=597$ Python, $n=807$ Node/TS.}
\label{tab:sink-families}
\small
\begin{tabular}{@{}lrr@{}}
\toprule
\textbf{Tracer family} & \textbf{Python \%} & \textbf{Node/TS \%} \\
\midrule
\multicolumn{3}{@{}l}{\emph{Low-specificity probes}$^{\dagger}$} \\
LogInjection             & 85.4 & 66.7 \\
SensitiveDataLogging     & 72.5 & 79.1 \\
FileAccess               & 59.5 & 50.3 \\
FileWrite                & 34.8 & 47.6 \\
ReDoS                    & 31.0 & 20.0 \\
BlockingEventLoopCall    & 30.2 & 39.7 \\
LoopCondition            & 27.8 & 16.9 \\
UnrestrictedIPBinding    & 26.0 & 37.1 \\
WeakPRNG                 & 16.4 & 24.0 \\
\midrule
\multicolumn{3}{@{}l}{\emph{Specific probes}} \\
NetworkExecution         & 57.8 & 61.5 \\
CommandExecution         & 30.8 & 33.0 \\
CodeExecution            & 24.3 & 24.7 \\
TempFileLeak             & 21.8 & 17.0 \\
SQLExecution             & 20.3 & 14.7 \\
UnverifiedCodeDownload   & 17.3 & 15.4 \\
DataAmplification        & 10.9 & 12.8 \\
WeakHash                 &  9.0 &  7.6 \\
XXE                      &  5.0 &  2.4 \\
SSLCertValidation        &  4.2 &  2.5 \\
Deserialization          &  4.0 &  4.6 \\
NonceReuse               &  3.2 &  7.1 \\
\midrule
matches no family        &  2.0 &  3.2 \\
\bottomrule
\end{tabular}

\vspace{0.4em}
{\footnotesize $^{\dagger}$~These probes match common language constructs
rather than distinctive sink APIs---\texttt{print(} for
\textsc{LogInjection}, \texttt{os.environ}/\texttt{process.env} for
\textsc{SensitiveDataLogging}, any \texttt{re.compile} for \textsc{ReDoS},
any \texttt{listen(} for \textsc{UnrestrictedIPBinding}. Their high rates
indicate that the tracer is broadly applicable, not that the behaviour is
widespread, and we do not use them for stratification. Appendix~\ref{app:probes} lists their patterns.}
\end{table}
